\section{De la semilla al corpus multilingüe explotable}

Un conjunto de evaluación responde «qué tan bien se hace»; los datos de entrenamiento además deben responder «de dónde se aprende». Para los idiomas de bajos recursos, la mayor dificultad no es la falta de un modelo lo bastante grande, sino que incluso falta un punto de partida confiable para language identification, vectores de oraciones, limpieza del corpus y backtranslation. NLLB-Seed y el mining pipeline posterior usan precisamente una pequeña cantidad de evidencia de alta calidad para poner en marcha una línea de producción de supervisión débil a mayor escala.

\subsection{Por qué no se puede empezar desde cero}

Esta sección pasa de la infraestructura de evaluación al problema de arranque de los datos de entrenamiento. La minería de datos suele describirse como «buscar oraciones paralelas en internet», pero para determinar a qué idioma pertenece una página web, si dos oraciones son semánticamente equivalentes o si un filtro descarta por error datos válidos, se necesitan anotaciones o representaciones confiables previas. La función de los seed data no es aportar la escala final, sino calibrar esos componentes básicos.

\lecturefigure{slide-14-seed-data-why.jpg}{La identificación de idioma, los vectores de oraciones y los modelos de traducción requieren datos semilla confiables}{Stanford, video de la clase 00:13:51}

\begin{knowledgebox}{Lectura de la figura: un mismo conjunto de seed data sirve a cuatro arrancadores}
Un texto confiable en el idioma de destino puede entrenar language identification; los pares de oraciones alineadas pueden entrenar un sentence encoder; el lado monolingüe puede apoyar el language modeling; el lado bilingüe puede poner en marcha un translation model. Si los datos iniciales están contaminados con idiomas equivocados, traducción automática o mezcla de escrituras, el error se propaga simultáneamente por las cuatro rutas.
\end{knowledgebox}

\begin{importantbox}{El efecto de palanca de los seed data}
El valor de una pequeña cantidad de datos humanos no depende solo del número de muestras, sino de cuántos componentes posteriores puede calibrar:
\[
\text{Seed}\rightarrow\{\text{LID},\text{Encoder},\text{Filter},\text{MT}\}
\rightarrow\text{Mined/BT Data}\rightarrow\text{Better MT}.
\]
Este es un bootstrapping loop. La calidad inicial determina el límite superior del error de las pseudoetiquetas y del corpus obtenido por minería en etapas posteriores; por ello, el valor marginal de la semilla puede ser mayor que el de la misma cantidad de texto web aleatorio.
\end{importantbox}

\subsection{NLLB-Seed: un recurso de arranque pequeño y confiable}

La slide de la clase informa que NLLB-Seed cubre 43 idiomas y alrededor de 6,193 oraciones, y explica que el contenido proviene de la lista de temas de cultura general de Wikimedia. Las cifras de la publicación final del artículo de NLLB varían ligeramente según la versión y la selección de idiomas; esta clase conserva las cifras de la slide y las interpreta como un recurso semilla de alta calidad de «unos miles de oraciones y decenas de idiomas», no como el grueso del entrenamiento a gran escala.

\lecturefigure{slide-15-nllb-seed.jpg}{Número de idiomas, número de oraciones y origen en temas de cultura general de NLLB-Seed}{Stanford, video de la clase 00:14:18}

\begin{knowledgebox}{Lectura de la figura: por qué se eligieron temas de cultura general y no páginas web aleatorias}
Los temas de cultura general abarcan personas, historia, geografía, ciencia y otros tipos de contenido, y reducen el ruido de plantillas de un solo sitio; además, las oraciones consecutivas dan contexto a los traductores. El costo es que el contenido de origen sigue la estructura de conocimiento de la Wikipedia en inglés y el texto de destino puede presentar translationese. El objetivo de la semilla es un arranque confiable, no representar todas las culturas locales ni los contextos de la lengua hablada.
\end{knowledgebox}

\teachervoice{Fan resume las tareas de generación de bajos recursos con «can't start from nothing»: sin unas cuantas oraciones confiables, es difícil saber en qué idioma está un texto y también construir la traducción inicial. El punto pedagógico de los seed data es la calidad del punto de partida, no una competencia por el número de muestras.}

\subsection{WikiMatrix: el antecedente de la minería a gran escala}

WikiMatrix mostró una posibilidad fundamental: usar vectores de oraciones multilingües para buscar oraciones semánticamente correspondientes en Wikipedia permite producir datos paralelos para un gran número de pares de idiomas. Para NLLB demostró que la embedding-based mining puede ampliar la cobertura, pero la escala y los temas de Wikipedia son limitados; para llegar realmente a doscientos idiomas hay que recurrir a la web abierta, mucho más desordenada.

\lecturefigure{slide-16-wikimatrix.jpg}{WikiMatrix usa vectores de oraciones multilingües para extraer oraciones paralelas a gran escala}{Stanford, video de la clase 00:15:00}

\begin{knowledgebox}{Lectura de la figura: los 135M y 1620 del título son escala, no garantía de calidad}
Cuantos más candidatos se extraen, más falsos positivos puede haber. La aportación de WikiMatrix fue convertir la recuperación por vectores de oraciones en un pipeline escalable; NLLB todavía necesitó un mejor encoder, identificación de idioma, margin score, filtrado y validación humana para convertir oraciones web similares en bitext utilizable para entrenamiento.
\end{knowledgebox}

\subsection{Sentence Alignment: de las páginas web a los pares de oraciones}

La clase divide la minería de pares de oraciones en tres niveles: buscar en Common Crawl contenido web posiblemente relacionado, extraer y limpiar el texto, y luego usar vectores de oraciones para encontrar coincidencias. La falla de cualquier nivel produce emparejamientos erróneos: las plantillas de las páginas pueden parecerse aunque el texto principal sea distinto, el segmentador de oraciones puede no adaptarse a la escritura y los vectores pueden acercar oraciones del mismo tema pero con hechos distintos.

\lecturefigure{slide-17-sentence-alignment.jpg}{Del emparejamiento de contenido de Common Crawl a la alignment a nivel de oración}{Stanford, video de la clase 00:15:39}

\begin{knowledgebox}{Lectura de la figura: la alignment no es coincidencia de cadenas}
El sistema primero reduce los documentos candidatos, luego codifica las oraciones en vectores compartidos y finalmente compara su cercanía semántica. Al leer la figura hay que distinguir content retrieval de sentence retrieval: la primera reduce el espacio de búsqueda; la segunda decide si se forma un ejemplo positivo. Si se aplica directamente el vecino más cercano global, las plantillas comunes y las oraciones cortas sepultan las traducciones verdaderas.
\end{knowledgebox}

Sean $f(x)$ y $g(y)$ los vectores de la oración de origen $x$ y de la oración de destino $y$, respectivamente; la similitud más básica es
\[
\operatorname{cos}(x,y)=\frac{f(x)^\top g(y)}{\lVert f(x)\rVert_2\lVert g(y)\rVert_2}.
\]
Aquí $f,g$ denotan el sentence encoder multilingüe; el numerador es el producto interno y el denominador normaliza los vectores. En la práctica, la mining suele usar un margin-based score, que compara la similitud del candidato con la similitud promedio de su vecindario local para reducir la interferencia de la hubness y de las oraciones genéricas.

\begin{lstlisting}[caption={Flujo conceptual de la minería de corpus paralelos con vectores de oraciones}]
for source_document in crawl:
    target_candidates = retrieve_related_documents(source_document)
    source_sentences = clean_and_segment(source_document)
    target_sentences = clean_and_segment(target_candidates)
    pairs = nearest_neighbors(source_sentences, target_sentences)
    keep = margin_filter(pairs, language_id=True, quality=True)
    send_selected_samples_to_bilingual_validation(keep)
\end{lstlisting}

\subsection{Encoder quality: el primer cuello de botella de los idiomas de bajos recursos}

Si el sentence encoder solo forma un buen espacio compartido para los idiomas de altos recursos, las oraciones de bajos recursos se agrupan con vecinos equivocados, y por grande que sea la mining posterior solo producirá ruido en masa. La clase compara masked language modeling con multilingual distillation: la primera aprovecha el contexto monolingüe; la segunda hace que un student multilingüe imite el espacio semántico de un teacher más fuerte.

\lecturefigure{slide-18-encoder-training.jpg}{Mejora del encoder con masked language modeling y multilingual distillation}{Stanford, video de la clase 00:16:33}

\begin{knowledgebox}{Lectura de la figura: los dos objetivos complementan, respectivamente, el modelado del lenguaje y la alineación entre idiomas}
El masked language modeling hace que el encoder aprenda la estructura interna del idioma de destino; la multilingual distillation usa oraciones semánticamente equivalentes para acercar distintos idiomas a un espacio compartido. La primera no garantiza por sí sola la alineación entre idiomas, y la segunda depende de un teacher confiable y de datos emparejados; NLLB combina ambas para aumentar el recall y la precision de la mining en idiomas de recursos extremadamente bajos.
\end{knowledgebox}

\begin{warningbox}{La puntuación del encoder se convierte en un sesgo de selección de datos}
El mining pipeline solo ve las similitudes que el encoder puede representar. Si cierta escritura, dominio o variante tiene mala calidad en el encoder, menos de sus oraciones entran al conjunto de entrenamiento y después al modelo le cuesta todavía más mejorar ese idioma, lo que forma una retroalimentación negativa del data flywheel.
\end{warningbox}

\subsection{LASER3: primero hacer más equitativa la representación, luego hablar de escala}

Ya se explicó que el encoder determina lo que la mining puede ver; esta sección verifica ese cuello de botella con los resultados de LASER3. La gráfica de resultados compara, por idioma, la calidad de los vectores de oraciones. Lo que la clase quiere que el lector observe no es que «todas las barras suben», sino que el nuevo encoder mejora notablemente muchos idiomas que antes eran muy débiles, lo que aumenta la probabilidad de que encuentren candidatos correctos en la minería web.

\lecturefigure{slide-19-laser3-encoder-quality.jpg}{Calidad del encoder interlingüe de LASER3 frente a las primeras versiones de LASER}{Stanford, video de la clase 00:18:48}

\begin{knowledgebox}{Lectura de la figura: primero buscar los idiomas con línea base baja y luego ver si la mejora es generalizada}
El eje horizontal son los idiomas y el eje vertical, una métrica de calidad relacionada con el encoder; las barras claras y oscuras corresponden a versiones distintas. Lo más importante no es el aumento promedio, sino si los idiomas de la cola pasan de ser casi inutilizables a poder sostener la mining. La figura tampoco demuestra que la calidad final de la traducción aumente en la misma proporción, porque después siguen los cuellos de botella de visibilidad web, filtrado y entrenamiento del modelo.
\end{knowledgebox}

\teachervoice{El docente llama a la encoder quality un low-resource challenge: un método de mining viable en idiomas de altos recursos, si la representación compartida falla para ciertos idiomas, los excluye antes de que lleguen al modelo.}

\subsection{Una fábrica de datos iterativa, no un ETL único}

La slide completa conecta en un ciclo language identification, monolingual cleaning, sentence mining, filtering, entrenamiento del modelo y bilingual validation. Cada mejora del modelo y del encoder puede redescubrir mejores datos; cada vez que la validación humana detecta un error sistemático, también hay que volver a la LID, a la limpieza o a la configuración de umbrales.

\lecturefigure{slide-20-data-pipeline.jpg}{Pipeline de datos de bajos recursos de NLLB y su ciclo de validación humana}{Stanford, video de la clase 00:21:15}

\begin{knowledgebox}{Lectura de la figura: seguir las flechas para encontrar la retroalimentación, no solo leer los recuadros}
Los LID training data sustentan la clasificación de idioma de las páginas web; los monolingual data limpios entran a la mining; el bitext candidato, tras el filtering, entrena el MT; las salidas del modelo y las muestras vuelven a ser validadas por personal bilingüe. Los múltiples ciclos de la figura significan que datos, modelo y evaluación iteran juntos, y que cualquier umbral fijo debe recalibrarse por idioma.
\end{knowledgebox}

\begin{importantbox}{Propagación del error en el pipeline}
Si la precision efectiva de la etapa $i$ es $p_i$ y las etapas se aproximan como filtros en serie, la confiabilidad aproximada de los candidatos finales está acotada por
\[
P_{\text{effective}}\approx\prod_i p_i
\]
Aquí $p_i$ representa la tasa de acierto de etapas como LID, limpieza, emparejamiento de documentos y puntuación de pares de oraciones. La expresión no es un modelo probabilístico exacto, pero muestra que varios componentes «bastante buenos» conectados en serie pueden perder calidad de forma significativa; por eso se necesitan validación humana por muestreo y corrección mediante ciclos.
\end{importantbox}

\subsection{Stopes: abrir el pipeline de investigación como infraestructura}

Esta sección pasa del flujo algorítmico a la ingeniería reproducible. A continuación, la clase presenta el punto de acceso de código abierto de Stopes. Lo que se abre no es solo el modelo final, sino también las herramientas para procesar páginas web monolingües, extraer bitext y reproducir la producción de datos. Este tipo de release permite que investigadores externos revisen la lógica de filtrado, sustituyan componentes de idioma y añadan datos locales, en lugar de limitarse a descargar un checkpoint inexplicable.

\lecturefigure{slide-21-stopes-open-source.jpg}{Stopes: la herramienta de minería de datos multilingües que NLLB publicó como código abierto}{Stanford, video de la clase 00:21:45}

\begin{knowledgebox}{Lectura de la figura: el nivel de apertura determina la profundidad de la reproducibilidad}
Abrir solo el modelo permite reproducir únicamente la inference; abrir el benchmark permite reproducir la evaluation; solo abrir las herramientas de LID, encoder, mining y filtering permite reconstruir los training data. Stopes está en el tercer nivel: no sustituye el conocimiento local de los idiomas, pero hace más verificables los supuestos de ingeniería del pipeline de datos.
\end{knowledgebox}

\subsection{El límite inferior estricto de los idiomas de recursos extremadamente bajos}

Por último, hay que delimitar el límite inferior estricto que el pipeline no puede superar y distinguir entre un recall algorítmico insuficiente y la inexistencia de suficiente texto en internet. Aun con mining a gran escala, algunos idiomas casi no tienen suficiente texto web monolingüe; una escritura particular puede hacer fallar la limpieza y la segmentación genéricas, y el contenido en línea puede concentrarse en lo religioso o en un solo dominio. En ese caso, «rastrear más páginas web» no produce automáticamente señales de entrenamiento diversas, equilibradas y confiables.

\lecturefigure{slide-22-very-low-resource-challenges.jpg}{Cuello de botella monolingüe, escritura y limitaciones de dominio en idiomas de recursos extremadamente bajos}{Stanford, video de la clase 00:22:36}

\begin{knowledgebox}{Lectura de la figura: tres desafíos corresponden a tres remedios no intercambiables}
La escasez de datos monolingües requiere datos de la comunidad, voz u otras fuentes; una escritura particular requiere normalization, segmentación y LID a la medida; la concentración en un dominio requiere ampliar deliberadamente los temas, no simplemente repetir el muestreo. Un modelo más grande no puede crear evidencia textual que no existe; solo ajusta con más detalle el sesgo existente.
\end{knowledgebox}

\begin{warningbox}{La escala no atraviesa el data floor}
Cuando la web presence se acerca a cero, el límite superior del recall de la mining lo determina la existencia de los datos; cuando las páginas web pertenecen a un solo dominio, aumentar la escala del crawl produce sobre todo repeticiones. Un plan de ingeniería debe diagnosticar por separado «faltan datos» y «el algoritmo no encontró los datos».
\end{warningbox}

\teachervoice{Fan afirma explícitamente que, para los idiomas de recursos extremadamente bajos, incluso la mining a gran escala es muy difícil, y que los monolingual data son el limiting factor. Este juicio reubica el problema: deja de ser la falta de capacidad de cómputo y pasa a ser la digitalización de los idiomas y la visibilidad de los datos.}

\subsection{Resumen de la sección}

NLLB usa una semilla de alta calidad para poner en marcha la LID, el encoder, la mining y el MT, y luego construye con LASER3 y Stopes una fábrica de datos iterativa. El límite superior del sistema está acotado a la vez por la calidad de la representación, la visibilidad web, el manejo de las escrituras, la diversidad de dominios y la validación humana.
